\chapter{Three Structural Conjectures}
\label{v4-arith:app:three-conjectures-dep}

The three conjectures F1, F2, F3 of Chapter~\ref{v4-arith:open-frontiers}
are hierarchically ordered: F3 depends on F1 and F2; F1 and F2 are
independent of each other but share local inputs from Fargues--Scholze
arXiv:2102.13459. The dependency graph, the local theorems each
conjecture imports, and the global condition carried by each
conjecture are recorded below.

\section{The Graph}
\label{v4-arith:app:C:graph}

\begin{center}
\begin{tikzcd}[row sep=large, column sep=large]
\text{F1: global arith.~AG} \ar[r] \ar[d] & \text{F2: perfectoid }E_{2} \ar[d] \\
\text{Local F1 (Fargues--Scholze)} \ar[ur, dashed] & \text{F3: arith.~Ben-Zvi--Nadler}
\end{tikzcd}
\end{center}

Each of F1 and F2 holds locally (the Fargues--Scholze and Emerton--Gee
inputs are theorems) but remains a conjecture globally. F3 combines
the local inputs of F1 and F2 with arithmetic Ben-Zvi--Nadler
machinery; the global identification of the Hochschild trace class
with \(\mathcal{V}^{\mathrm{prim}}\) is conditional on F1 and F2.

\section{F1: Global Arithmetic Arinkin--Gaitsgory}
\label{v4-arith:app:C:F1-details}

\begin{conjecture}[restatement of F1]
\label{v4-arith:app:C:F1-restatement}
For \(G=GL_{2}\) and the Selmer-derived global spectral stack
\(\mathcal{X}^{\mathrm{Sel}}_{\bar\rho,\det,\mathcal{L}}\) (the
derived moduli stack of Selmer deformations of a residual
representation \(\bar\rho\) with determinant condition \(\det\) and
cyclotomic twist datum \(\mathcal{L}\), as in
Chapter~\ref{v4-arith:open-frontiers}):
\begin{equation}
\mathrm{D}([GL_{2}])^{\mathrm{sph\text{-}fin}}\;\simeq\;\mathrm{IndCoh}_{\mathrm{Nilp}}\bigl(\mathcal{X}^{\mathrm{Sel}}_{\bar\rho,\det,\mathcal{L}}\bigr).
\end{equation}
\end{conjecture}

\paragraph{Local inputs.}
\begin{enumerate}
\item Fargues--Scholze arXiv:2102.13459 (local geometrization and
spectral action, not the displayed EG equivalence).
\item Caraiani--Scholze arXiv:1511.02418 (unitary Shimura variety case).
\item Emerton--Gee stack construction arXiv:1908.07185.
\end{enumerate}

\paragraph{Global condition.}
\begin{itemize}
\item Proof of the local automorphic-to-Emerton--Gee
\(\mathrm{IndCoh}_{\mathrm{Nilp}}\) comparison, including small primes
and central-character families.
\item Compatibility of local comparison functors under
restricted-tensor-product gluing over primes.
\item Identification of the Whittaker model across all places
simultaneously.
\item Functional-equation-level compatibility between spectral and
automorphic sides.
\end{itemize}

\paragraph{Hypothesis.} This is a global comparison theorem.

\section{F2: Perfectoid \(E_{2}\)-Enhancement}
\label{v4-arith:app:C:F2-details}

\begin{conjecture}[restatement of F2]
\label{v4-arith:app:C:F2-restatement}
The local Iwahori Hecke category \(\mathrm{Hk}^{\mathrm{Iw}}_{GL_{2},p}\)
carries a canonical \(E_{2}\)-monoidal structure from the
perfectoid Beilinson--Drinfeld Grassmannian over \(\mathrm{Spec}\,\mathbb{Z}_{p}\),
and the restricted tensor product
\(\bigotimes'_{p}\mathrm{Hk}^{\mathrm{Iw}}_{GL_{2},p}\) inherits a
globally-defined \(E_{2}\)-monoidal structure.
\end{conjecture}

\paragraph{Local ingredients (theorem).}
\begin{enumerate}
\item Perfectoid affine Grassmannian (Fargues--Scholze, Scholze
diamonds).
\item Witt-vector affine Grassmannian (Zhu arXiv:1603.05593;
Bhatt--Scholze arXiv:1507.06490).
\item Equal-characteristic BD Grassmannian \(E_{2}\)-structure
(Mirkovi\'c--Vilonen; Beilinson--Drinfeld Grassmannian factorization).
\end{enumerate}

\paragraph{Global condition.}
\begin{itemize}
\item Diamond-gluing of Iwahori Hecke categories across primes.
\item Compatibility with Satake and Whittaker equivalences.
\item Perfectoid Dunn-additivity argument producing \(E_{1}\)-algebra
structure on the Hochschild trace.
\end{itemize}

\paragraph{Hypothesis.} This is a mixed-characteristic factorisation theorem.

\section{F3: Arithmetic Ben-Zvi--Nadler}
\label{v4-arith:app:C:F3-details}

\begin{conjecture}[restatement of F3]
\label{v4-arith:app:C:F3-restatement}
\(\mathrm{HH}_{\bullet}\bigl(\mathrm{Hk}^{\mathrm{Iw}}_{GL_{2},\mathrm{glob}}\bigr)\xrightarrow{\sim}\mathcal{V}^{\mathrm{prim}}\)
as \(E_{1}\)-algebras.
\end{conjecture}

\paragraph{Ingredients.}
\begin{enumerate}
\item Conjecture F1 (for the spectral side identification).
\item Conjecture F2 (for \(E_{2}\)-enhancement of the local Hecke
factors).
\item Ben-Zvi--Nadler arXiv:0904.1247 (for the \(\mathbb{C}\)-version,
which F3 arithmeticises).
\item Explicit computation matching the Hochschild trace character
to \(\zeta(s)\).
\end{enumerate}

\paragraph{Consequences of F3.}
\begin{itemize}
\item F3 implies the Gaiotto Chern--Simons identification
(\Cref{v4-arith:thm:arith-CS-VA-correspondence}).
\item F3 implies the Bezrukavnikov--Bernstein Hochschild identification
(\Cref{v4-arith:thm:crowning-placement}).
\item F3 implies Beilinson's bar--cobar identification
(\Cref{v4-arith:thm:bar-verifications}).
\item F3 implies the arithmetic six-window placement
(\Cref{v4-arith:thm:six-window-placement}).
\end{itemize}

\paragraph{Hypothesis.} This is an arithmetic trace theorem.

\section{Independent Conjecture: The \(\mathsf{G}\)-Row Complementarity}
\label{v4-arith:app:C:G-row-independent}

\begin{conjecture}[\(\kappa+\kappa^{!}=0\Leftrightarrow\xi(s)=\xi(1-s)\)]
\label{v4-arith:app:C:G-row-statement}
For the arithmetic Heisenberg \(\mathcal{V}^{\mathrm{prim}}_{\mathrm{Heis}}\),
the Vol~I derived-centre complementarity
\(\kappa^{\mathrm{Ar}}+(\kappa^{!})^{\mathrm{Ar}}=0\) is equivalent
to the Riemann functional equation
\(\xi(s)=\xi(1-s)\).
\end{conjecture}

This is \emph{independent} of F1, F2, F3: its proof proceeds
directly from the arithmetic Koszul-resolution computation of
Hochschild cochains on the arithmetic Heisenberg, using Tate's
adelic Poisson summation. A proof of this conjecture would state
the arithmetic \(\mathsf{G}\)-row as a Vol~I Theorem~C corollary.

\paragraph{Hypothesis.} This is an arithmetic Koszul-duality theorem.

\section{Chain of Implications}
\label{v4-arith:app:C:chain-of-implications}

The arithmetic branch has the following implication form:

\begin{enumerate}
\item The \(\mathsf{G}\)-row conjecture
(\Cref{v4-arith:app:C:G-row-statement}) would identify the Riemann
functional equation as an arithmetic Koszul-duality consequence.
\item F3 conditional on F1 and F2 unifies the three presentations of
\(\mathcal V^{\mathrm{prim}}\).
\item F1 is the global arithmetic Arinkin--Gaitsgory input.
\item F2 is the perfectoid \(E_{2}\)-enhancement input.
\item F1, F2, and the conditional F3 comparison imply the six-window
structural statement.
\end{enumerate}

\section{Connection to Vol~III}
\label{v4-arith:app:C:vol3-connection}

Vol~III's \(\mathsf{B}\)-row \(\kappa+\kappa^{!}=8\) (Mukai-enhanced K3
Heisenberg via Bruinier Heegner Chern-class reciprocity) has an
arithmetic avatar (F4, Chapter~\ref{v4-arith:open-frontiers}~\S3).
Together the arithmetic \(\mathsf{G}\)-row and the arithmetic
\(\mathsf{B}\)-row would give two of the five Vol~I
derived-centre-ceiling values as consequences of automorphic and
Galois-deformation theorems.  The chiral centre is then constrained by
arithmetic comparison data.

\section{Extended Graph with F-RH}
\label{v4-arith:app:C:extended-graph}

The conjectures F1, F2, F3 of the preceding sections address the
six-window placement of \(\mathcal{V}^{\mathrm{prim}}\). A second
cluster of conditions, independent of F1--F3, governs the
connection to the Riemann hypothesis. Chapter~\ref{v4-arith:closure}
provides the structural branch; Chapter~\ref{v4-arith:rh-reformulation}
states F-RH, whose zero-placement clause is RH and whose other clauses
are the positive Hermitian, Koszul, and Zhu structures on
\(\mathcal{V}^{\mathrm{prim}}\). The arrows are:

\paragraph{F1--F3 structural branch (six-window placement).}
\begin{center}
\begin{tikzcd}[row sep=small, column sep=small]
\text{F1.a (local comparison, TW)} \ar[dr] & \text{F1.c (Selmer target)} \ar[d] & \text{F2.a (spherical input)} \ar[d] \\
\text{L}_{\mathrm{ACD}} \ar[r] & \text{F1 conditional} \ar[r] & \text{F2 conditional} \ar[d] \\
\text{F2.c (\(\infty\)-Tate)} \ar[ur] & \text{F3.BZN}^{\mathrm{Ar}} \ar[r] & \text{F3 conditional}
\end{tikzcd}
\end{center}

\paragraph{four-input RH branch.}
\begin{center}
\begin{tikzcd}[row sep=small, column sep=small]
\text{arith.~Positselski} \ar[r] & \text{P1 functional equation} \ar[dr] & & & \\
\text{Tate 1950} \ar[r] & \text{P2 (Tate)} \ar[r] & \mathrm{H}\text{-}\mathrm{P}^{\mathrm{Ar}}+\mathrm{VB}^{*} \ar[r] & \text{zero-placement} \ar[r, "\Leftrightarrow"] & \text{RH} \\
\text{Koszul--Pet compat} \ar[r] & \text{P3 restricted unitarity} \ar[ur] & & &
\end{tikzcd}
\end{center}

The dependency partitions into two branches:

\paragraph{F1--F3 structural branch (six-window placement, RH-independent).}
\(\mathrm{L}_{\mathrm{ACD}}\) (the adelic comparison and descent
lemma relating the global spherical category to the restricted tensor
product of local Iwahori categories) \(\to\) F1 \(\to\)
F3.BZN\(^{\mathrm{Ar}}\) (the arithmetic Ben-Zvi--Nadler functoriality
for the loop space of the global spectral stack) \(\to\) F3. The
sub-conjecture F2.c has a restricted pointed central-unit theorem in
Chapter~\ref{v4-arith:tate-tensor}; the mixed-characteristic Iwahori
\(E_{2}\) local input (F2.b) remains conditional.  The F1--F3
structural branch states the six conditional windows of
\(\mathcal{V}^{\mathrm{prim}}\).

\paragraph{four-input RH branch.}
Here P1 is the functional-equation branch (the statement that the
adelic Pozitsel\textquotesingle ski bar complex has the
functional-equation symmetry of \(\xi(s)\)); P2 is the
Tate-adelic computation (Tate 1950);
P3 is restricted Koszul unitarity (the Koszul--Petersson compatibility
on the spherically-finite cuspidal core); and H-P\(^{\mathrm{Ar}}\)
is the arithmetic Hilbert--P\'olya determinant-trace problem
(existence of a self-adjoint operator whose normal translate has
regularised determinant divisor \(\mathrm{div}\,\xi_{\mathbb R}\)).
The implication chain is: arithmetic Positselski
\(\to\) P1; Koszul--Pet compatibility \(\to\) P3; (P1, P2, P3,
H-P\(^{\mathrm{Ar}}\)) \(\to\) zero-placement under VB\(^{*}\). Here
VB\(^{*}\) is the arithmetic Virasoro--Bogomolov positivity input:
the zero-mode positivity of the Hamburger spectral measure,
equivalent to RH. Under VB\(^{*}\), the four inputs force
zero-placement, hence RH.

The two branches are logically independent.  The F1--F3 structural
branch constructs the trace object.  The four-input RH branch controls
zero placement.  A theorem on one branch is not a theorem on the other.

\section{Boundary Conditions}
\label{v4-arith:app:C:residual-list}

\begin{table}[h]
\centering
\small
\begin{tabular}{l|l|l|l}
Input & Branch & Size & Present form \\
\hline
\(\mathrm{L}_{\mathrm{ACD}}\) (adelic descent) & A & global theorem & Conjectural \\
F2.c (\(\infty\)-Tate product) & A & mixed-characteristic theorem & Restricted pointed central-unit theorem; full mixed-characteristic statement conditional \\
F2.b (mixed-char Iwahori \(E_{2}\)) & A & mixed-characteristic theorem & Conditional \\
F3.BZN\(^{\mathrm{Ar}}\) (arith.~BZN) & A & trace theorem & Conjectural \\
arith.~Positselski & B & Koszul theorem & Strict abelian Heis. model constructed; full domain separate \\
Koszul--Pet compat & B & automorphic theorem & Sph-fin cuspidal core theorem \\
VB\(^{*}\) (arith.~Virasoro positivity) & B & Unknown & Positive spectral condition implies RH \\
\end{tabular}
\caption{Boundary conditions for the structural and RH branches.}
\label{v4-arith:app:C:residual-table}
\end{table}

\section{A Sufficient Path to RH}
\label{v4-arith:app:C:sufficient-path}

The sufficient implication to RH inside the arithmetic formalism uses
four inputs from the four-input RH branch: the strict abelian-Heisenberg arithmetic
Positselski bar, Koszul--Pet compatibility on the normalised core,
H-P\(^{\mathrm{Ar}}\), and the positive VB\(^{*}\) spectral condition. The
text proves the restricted pointed central-unit form of F2.c, constructs
the arithmetic Positselski bar as a strict reciprocity-closed chain complex, and
Koszul--Pet is established on the sph-fin cuspidal core.
Determinant-trace realisation, full-sector P3,
mixed-characteristic Iwahori \(E_{2}\), and
VB\(^{*}\) positivity remain separate conditions.

\subsection{Refined Boundary List}
\label{v4-arith:app:C:residual-list-w2}

\begin{table}[h]
\centering
\small
\begin{tabular}{l|l|l|l}
Input & Branch & Size & Present form \\
\hline
\(\mathrm{L}_{\mathrm{ACD}}.0,.1,.2,.3\) & A & each large theorem & Partial reductions; global glue remains conditional \\
(BC-diamond-glue).core & A & global theorem & Conjectural, single cross-prime Bernstein compat \\
F2.c (pointed \(E_{n}\) datum) & A & -- & Restricted theorem of \Cref{v4-arith:tate:thm:F2c-main}; mixed-characteristic Iwahori \(E_2\) input separate \\
F3.BZN\(^{\mathrm{Ar}}\) R1.fin-type & A & -- & External finite-type input; arithmetic application conditional \\
F3.BZN\(^{\mathrm{Ar}}\) R1.formal.a & A & -- & Trace theorem after the categorical colimit of \(\mathrm{IndCoh}_{\mathrm{Nilp}}\) is known \\
F3.BZN\(^{\mathrm{Ar}}\) R1.formal.b,c & A & pro-\'etale theorem & Pro-\'etale loop-space and flat-locus hypotheses \\
F3.BZN\(^{\mathrm{Ar}}\) R2 (Selmer target) & A & Selmer theorem & Target and tangent complex fixed; Noetherian and nilpotent-support hypotheses remain \\
F3.BZN\(^{\mathrm{Ar}}\) R3 (full target) & A & -- & Conditional; the inertia shadow is not the Weil/\((\varphi,\Gamma)\) target \\
F3.BZN\(^{\mathrm{Ar}}\) R4 wild \(p\in\{2,3\}\) & A & -- & Wild summand included locally; global wild functoriality open \\
Wild Langlands functoriality & A & -- & Conjectural (\Cref{v4-arith:bzn-rcompletion:conj:wild-langlands-functorial}) \\
arith.~Positselski C1--C5 & B & -- & Strict completed model constructed for the abelian primary bar \\
arith.~Positselski (abelian Heis.) & B & -- & Forward Tate functional equation; converse under Positselski hypotheses \\
Koszul--Pet K5 (Godement--Jacquet) & B & -- & Conditional normalisation on the sph-fin cuspidal core \\
Koszul--Pet compat & B & -- & Sph-fin cuspidal core theorem; full-sector extension separate \\
H-P\(^{\mathrm{Ar}}\) & B & -- & Friedrichs domain theorem; determinant divisor and trace matching open \\
VB\(^{*}\) \(\Rightarrow\) zero-placement & B & -- & Conditional implication from P2, P3, H-P\(^{\mathrm{Ar}}\), and positive trace \\
VB\(^{*}\) (main direction) & B & -- & Positive spectral condition; determinacy alone is insufficient \\
VB\(^{*}\) converse & B & -- & Strictly finer than RH; 4-way equivalence (Ch.~\ref{v4-arith:HPAr-vbconverse}) \\
\end{tabular}
\caption{Refined boundary conditions.}
\label{v4-arith:app:C:residual-table-w2}
\end{table}

\subsection{Sufficient RH Path}

A sufficient implication to RH within the arithmetic formalism is the
four-input conditional conjunction on the four-input RH branch:

\begin{enumerate}
\item \textbf{arith.~Positselski} (the four-input RH branch): strict completed
Gersten--Parshin bar model on the abelian primary Heisenberg domain.
\item \textbf{Koszul--Pet compat} (the four-input RH branch): stable on the normalized
sph-fin cuspidal core; full \(V^{\mathrm{prim}}\) compatibility and
KMS\(_1\)/GNS identification remain open.
\item \textbf{H-P\(^{\mathrm{Ar}}\)} (the four-input RH branch): the independent
Hilbert--P\'olya determinant-divisor matching problem.
\item \textbf{VB\(^{*}\) main direction} (the four-input RH branch): the positive
Hamburger zero-mode spectral condition.
\end{enumerate}

The F1--F3 structural branch remains RH-independent: (BC-diamond-glue).core,
F2 trace continuity, and arithmetic BZN formal-loop compatibility affect the
construction of \(\mathcal{V}^{\mathrm{prim}}\), not zero
placement. VB\(^{*}\) remains RH-equivalent; no formalism shortens it,
but the ladder of equivalent statements (Li positivity, Weil
positivity, Hamburger determinacy, Hermite--Biehler, B\'aez-Duarte
sharp rate) is listed in
Chapter~\ref{v4-arith:HPAr-vbconverse}.

\section{Dependency Boundary}
\label{v4-arith:app:C:summary}

Three hierarchically ordered conjectures (F1 \(\to\) F2 \(\to\) F3)
and one independent conjecture (\(\mathsf{G}\)-row) govern the
synthesis branch. On the RH branch: P1 (functional equation),
P2 (Tate adelic), P3 (restricted Koszul unitarity),
H-P\(^{\mathrm{Ar}}\) (Hilbert--P\'olya determinant-divisor matching), and
VB\(^{*}\) (arithmetic Virasoro--Bogomolov zero-mode positivity,
RH-equivalent). The restricted pointed central-unit theorem gives the
proved part of F2.c, and the strict abelian primary arithmetic
Positselski bar supplies the P1 forward direction. The arithmetic branch as stated
(\Cref{v4-arith:closure}, \Cref{v4-arith:rh-reformulation}) reduces
the RH-strength condition to H-P\(^{\mathrm{Ar}}\) determinant-trace
realisation and the positive Hamburger VB\(^{*}\) condition.
